\section{Limitations}
\label{sec:limitations}

HARP is a controlled study of harm amplification, not a deployment guarantee. Its finance-oriented seven-agent setting uses synthetic state, simulated tools, resettable databases, and a deterministic decision gate, enabling auditable traces and reproducible comparisons but limiting external validity. Other domains, architectures, and escalation policies may amplify or suppress perturbations differently. The harness covers bounded role perturbations, including redirects, numeric manipulation, objective rewriting, shared-context corruption, collusion, and temporal triggers, but excludes training-time compromise, backdoors, malicious tools, network attacks, code tampering, side channels, and adaptive attackers. HARP needs trace visibility; broader models and domain-specific harm metrics remain needed.

\section{Conclusion}
Multi-agent LLM systems can fail in ways that are invisible to single-agent attack-success metrics. A locally plausible specialist deviation, shared biased evidence, or delayed memory trigger can be amplified by routing, aggregation, and final decision gates into high-impact downstream harm. \proposed~{} introduces a trace-first benchmark and amplification metric for measuring this phenomenon in a finance-oriented multi-agent architecture. The experiments show that prompt-only defenses preserve utility but leave amplification largely intact, tool-level guards trade security for utility and latency, and post-hoc trace integrity checking substantially reduces attack success while introducing false positives. These results argue for evaluation protocols that record full traces, compare clean and perturbed executions, and measure the local-to-global conversion of deviation into harm.
